% Printed 443.
\par\clearpage
\noindent\null\hfill{\small\textbf{HÖHENMESSUNGEN.}}\hfill\makebox[0pt][r]{\fontsize{13}{16}\selectfont 443}\par\medskip
\begin{center}
[3].\par\medskip
[Der Refraktionskoeffizient aus den Höhenmessungen bei der Gradmessung\par
und ihrer Fortsetzung bis Jever.]\par\bigskip
\textsc{Gauss} an \textsc{Schumacher}.\quad Göttingen, 27.~Dezember 1846.
\end{center}
\par\noindent\hspace*{\parindent}\textsc{Bessel} erwähnt [in einem hinterlassenen Aufsatze] meines Resultats für
die irdische Strahlenbrechung und meint, die Ursache, warum ich den Koeffizienten kleiner finde, als andere Astronomen, sei, weil ich nur an sonnigen
Tagen beobachtet habe, die anderen aber ohne Unterschied an sonnigen und
bedeckten. Diese Äußerung ist in vielfacher Beziehung nicht richtig. Erstlich ist das Faktische nicht unbedingt gültig. Vermutlich ist jener Aufsatz
früher geschrieben, als das Buch über die preußische Gradmessung; in letzterer, S.~197, hat \textsc{Bessel} auch andere Resultate angeführt, \textls[120]{zwischen} denen
meines, wie er es angibt, liegt. Dann hat \textsc{Bessel} bloß das Resultat aus
meinen Messungen von 1821 und 1822 gekannt, wie es im Jahrbuch für 1826
gedruckt ist. Die späteren Jahre geben fortschreitend größere Koeffizienten,
nämlich
\begin{center}
\begin{tabular}{r@{\ .\ .\ .\ }l}
1823&\(0{,}14125\)\\
1824&\(0{,}14778\)\\
1825&\(0{,}15826\).
\end{tabular}
\end{center}
\par\noindent\hspace*{\parindent}Das Mittelresultat aus allen 5 Jahren und zwar aus einer Anzahl von
Dreiecksseiten, deren Summe 290 geogr. Meilen beträgt, ist \(0{,}13974\), also
größer, als das Resultat aus \textsc{Bessels} eigenen Beobachtungen. Jenes Fortschreiten kann ich aber recht gut erklären. Es ist ungenau, zu sagen, dass
an sonnigen Tagen die Refraktion kleiner sei als an bedeckten. Die [Ursache]
ist vielmehr die: an sonnigen Tagen hängt die Refraktion in hohem Grade \textls[120]{von
der Tageszeit} ab, ohne allen Vergleich mehr, als an bedeckten. Mittags und
in den dem Mittag nächstgelegenen Vormittags- und Nachmittagsstunden ist die
Refraktion am kleinsten, in den späteren Nachmittagsstunden nimmt sie fortwährend mit einer außerordentlichen Regelmäßigkeit zu und gegen Sonnenuntergang ist sie gewiss nicht kleiner, sondern eher größer, als an bedeckten
Tagen. Die Zeit der kleinsten Refraktion ist zugleich die, wo die Luft feinen
\par\medskip\noindent\null\hfill\makebox[0.18\linewidth][l]{56\textsuperscript{*}}\par

% Printed 444.
\par\clearpage
\noindent\makebox[0pt][l]{\fontsize{13}{16}\selectfont 444}\hfill{\small\textbf{BRIEFWECHSEL.\quad HÖHENMESSUNGEN.}}\hfill\null\par\medskip
\noindent Beobachtungen am ungünstigsten ist. Da ich die Höhenmessungen nur wie
ein sekundäres Geschäft betrachtete, so verwandte ich darauf in den ersten
Jahren vorzugsweise die den feineren Horizontalwinkelmessungen ungünstigsten
Stunden, namentlich auch vormittags und mittags. In den folgenden Jahren
wurde die Zuziehung der Vormittagsstunden immer mehr beschränkt und hörte
zuletzt fast ganz auf, so dass häufig die Zenitdistanzen auch in ziemlich
späten Nachmittagsstunden gemessen wurden. Den von \textsc{Bessel} angeführten
Grund, woraus er die größeren von den französischen Astronomen gefundenen
Resultate erklären will, nämlich weil die \textsc{Borda}schen Kreise die Zenitdistanzen zu klein gäben, halte ich für sehr unerheblich. Es könnte daraus nur
eine sehr geringe Wirkung erfolgen; auch sind meine eigenen Messungen
ebenfalls mit einem \textsc{Borda}schen Kreise gemacht.
\begin{center}\rule{.16\linewidth}{.4pt}\end{center}
\begin{center}\textbf{BEMERKUNGEN.}\end{center}
\begingroup\small
\par\noindent\hspace*{\parindent}Zu dem ersten Teil des Aufsatzes aus \textsc{Bodes} astronomischem Jahrbuch sei bemerkt, dass die endgültigen Werte der geographischen Positionen bereits auf S.~427/428 gegeben sind. Der Abdruck des
letzten der beiden vorstehenden Briefe an \textsc{Schumacher} ist nach dem Original erfolgt, für den ersten ist
eine nach dem Original angefertigte Kopie benutzt worden. Die in ihm enthaltenen Zahlenwerte konnten
durch die Angaben eines \textsc{Gauss}schen Handbuches kontrolliert werden. Nach diesem sind die in der Tabelle,
S.~439, angegebenen Werte bei der Linie Lichtenberg-Deister der Reihe nach durch die folgenden zu ersetzen: \(1566{,}9\), \(222{,}3\), \(0{,}1419\). Die Höhenbeobachtungen während der 5 Jahre der Gradmessung haben (nach
verschiedenen Stellen des \textsc{Gauss}schen Handbuches) zur Ableitung des Refraktionskoeffizienten folgende Resultate ergeben:
\begin{center}
\setlength{\tabcolsep}{4pt}
\begin{tabular}{|l|r|r|r|r|}\hline
\multicolumn{1}{|c|}{Jahr der Beobachtungen}&\multicolumn{1}{c|}{\shortstack{\rule{0pt}{1pt}\\Anzahl\\der Linien}}&\multicolumn{1}{c|}{\shortstack{\rule{0pt}{1pt}\\Summe der\\Krümmungen}}&\multicolumn{1}{c|}{\shortstack{\rule{0pt}{1pt}\\Summe der\\Refraktionen}}&\multicolumn{1}{c|}{\shortstack{\rule{0pt}{1pt}\\Refraktions-\\koeffizient}}\\\hline
1821 u. 1822 (bis Wulfsode-Wilsede)&28&\(33203^{\prime\prime}{,}272\)&\(4309^{\prime\prime}{,}986\)&\(0{,}12981\)\\
1823 (bis Wilsede-Nindorf)&4&\(1778{,}089\)&\(178{,}572\)&\(0{,}12643\)\\
\phantom{1823 }(bis Hamburg)&5&\(6105{,}219\)&\(862{,}415\)&\(0{,}14126\)\\
1824 (bis Brillit-Garlste)&22&\(22549{,}844\)&\(3332{,}329\)&\(0{,}14778\)\\
1825 (bis Varel-Jever)&8&\(7763{,}012\)&\(1228{,}541\)&\(0{,}15826\)\\\hline
\multicolumn{1}{|r|}{Summe}&67&\(71399{,}436\)&\(9911{,}843\)&\(0{,}13882\)
\end{tabular}
\end{center}
\par\noindent\hspace*{\parindent}\textsc{Gauss} scheint die Ergebnisse der ersten 4 Linien von 1823 bei der Bildung des Mittelresultats,
S.~443, ausgeschlossen zu haben; in diesem Falle wird der Refraktionskoeffizient: \(0{,}13980\).
\par\noindent\hfill\textsc{Krüger}.\par
\endgroup
\begin{center}\rule{.2\linewidth}{.4pt}\end{center}

% Printed 445.
\par\clearpage
\noindent\null\hfill{\small\textbf{NACHLASS.\quad HÖHENMESSUNGEN.}}\hfill\makebox[0pt][r]{\fontsize{13}{16}\selectfont 445}\par
\vspace*{2.4cm}
\begin{center}{\Large\textbf{\textls[140]{NACHLASS.}}}\par\vspace{1.3cm}
Terrestrische Refraktion.\par\medskip
[1.]
\end{center}
\par\noindent\hspace*{\parindent}Es wird angenommen, dass der Weg des Lichtstrahls in einer einzigen Ebene
bleibt, und dass die Richtungen der Schwere an den verschiedenen Punkten
jenes Weges in einem einzigen Punkt \(C\) zusammentreffen; ebenso, dass die Veränderungen der Richtung des Weges den Veränderungen des Winkels an \(C\)
proportional sind.
\par\noindent\hspace*{\parindent}Es seien \(P\) und \(P'\) zwei [aufeinanderfolgende] Punkte der Lichtkurve;
\(90^\circ+u\) der Winkel zwischen \(PP'\) und \(PC\); \(PC=r\); \(A\) der Anfangspunkt
jener Linie, \(B\) der Endpunkt; \(\theta\) der Winkel zwischen \(CP\) und \(CA\); \(n\mathrm{d}\theta\) die
Richtungsveränderung; endlich seien \(a\) und \(b\) die Werte von \(r\) in \(A\) und \(B\),
und \(\alpha\) und \(\beta\) die Werte von \(u\) in diesen Punkten, d.~i. \(\alpha\) die Höhe, in welcher
\(B\) in \(A\) erscheint, \(\beta\) die Vertiefung, in welcher \(A\) in \(B\) erscheint.
\par\noindent\hspace*{\parindent}Man hat
\[u=\alpha+(1-n)\theta\]
\[\mathrm{d}r=r\mathrm{d}\theta.\operatorname{tang}u=r\mathrm{d}\theta.\operatorname{tang}(\alpha+(1-n)\theta)\]
\noindent [oder]
\[\log r=-\tfrac{1}{1-n}\log\cos(\alpha+(1-n)\theta)+\mathrm{const.}\]
\[\tfrac{\mathrm{const.}}{r^{1-n}}=\cos(\alpha+(1-n)\theta).\]
\par\noindent\hspace*{\parindent}Folglich [da für \(r=a,\quad\theta=0\) ist]
\[\tfrac{a^{1-n}\cos\alpha}{r^{1-n}}=\cos(\alpha+(1-n)\theta).\]

% Printed 446.
\par\clearpage
\noindent\makebox[0pt][l]{\fontsize{13}{16}\selectfont 446}\hfill{\small\textbf{NACHLASS.}}\hfill\null\par\medskip
\begingroup
\setlength{\abovedisplayskip}{4pt}\setlength{\belowdisplayskip}{4pt}
\setlength{\abovedisplayshortskip}{2pt}\setlength{\belowdisplayshortskip}{2pt}
\par\noindent\hspace*{\parindent}Ist also in \(B\)
\[\theta=t\quad[\text{und }r=b],\]
\noindent so ist
\[\beta=\alpha+(1-n)t\]
\noindent [und]
\[\tfrac{a^{1-n}\cos\alpha}{b^{1-n}}=\cos\beta,\]
\noindent [also wird]
\[\tfrac{\cos\alpha-\cos\beta}{\cos\alpha+\cos\beta}=\operatorname{tang}\tfrac12(\beta-\alpha)\operatorname{tang}\tfrac12(\beta+\alpha)=\tfrac{b^{1-n}-a^{1-n}}{b^{1-n}+a^{1-n}}.\]
\par\noindent\hspace*{\parindent}Man setze
\[\begin{aligned}\tfrac12(b+a)&=R+k\\b-a&=h,\end{aligned}\]
\noindent also
\[\begin{aligned}b&=R+k+\tfrac12h\\a&=R+k-\tfrac12h;\end{aligned}\]
\noindent [dann ist:]
\[\operatorname{tang}\tfrac12(\beta-\alpha)\operatorname{tang}\tfrac12(\beta+\alpha)=\frac{\left(1+\tfrac{h}{2(R+k)}\right)^{1-n}-\left(1-\tfrac{h}{2(R+k)}\right)^{1-n}}{\left(1+\tfrac{h}{2(R+k)}\right)^{1-n}+\left(1-\tfrac{h}{2(R+k)}\right)^{1-n}},\]
\noindent oder hinlänglich genau
\[=\tfrac{(1-n)h}{2(R+k)}.\]
\par\noindent\hspace*{\parindent}Die Länge des auf die Meeresfläche projizierten Bogens \(AB=\Delta\) gesetzt, wird
\[\Delta=Rt.\]
\par\noindent\hspace*{\parindent}Ferner ist
\[\text{I.}\hspace{3cm}\beta-\alpha=(1-n)t;\]
\noindent folglich
\[\begin{aligned}
(1-n)h&=2(R+k)\operatorname{tang}\tfrac12(\beta-\alpha)\operatorname{tang}\tfrac12(\beta+\alpha)\\
&=\tfrac{\operatorname{tang}\frac12(\beta-\alpha)}{\frac12(\beta-\alpha)}.(R+k)(\beta-\alpha)\operatorname{tang}\tfrac12(\beta+\alpha),\\
h&=\tfrac{\operatorname{tang}\frac12(\beta-\alpha)}{\frac12(\beta-\alpha)}.\tfrac{R+k}{R}.\Delta.\operatorname{tang}\tfrac12(\beta+\alpha),
\end{aligned}\]
\noindent wofür man auch
\[\text{II.}\hspace{3cm}h=\tfrac{R+k}{R}.\Delta.\operatorname{tang}\tfrac12(\beta+\alpha)\]
\noindent schreiben kann.
\par\noindent\hspace*{\parindent}Aus I findet man, wenn \(\alpha\) und \(\beta\) beide bekannt sind, \(n\); aus II \(h\).
\endgroup

% Printed 447.
\par\clearpage
\noindent\null\hfill{\small\textbf{HÖHENMESSUNGEN.}}\hfill\makebox[0pt][r]{\fontsize{13}{16}\selectfont 447}\par\medskip
\begingroup
\setlength{\abovedisplayskip}{5pt}\setlength{\belowdisplayskip}{5pt}
\begin{center}
[2.]\par\medskip
\begin{tabular}{l@{\ }l}
Depression des Horizonts\ .\ .\ .\ .\ .\ .\ .&\(\delta\)\\
Halbmesser der Erde\ .\ .\ .\ .\ .\ .\ .\ .\ .\ .&\(R\)\\
Höhe des Auges\ .\ .\ .\ .\ .\ .\ .\ .\ .\ .\ .\ .\ .\ .&\(h\)\\
Amplitude des Erdbogens\ .\ .\ .\ .\ .\ .&\(v\)\\
Ganze Refraktion\ .\ .\ .\ .\ .\ .\ .\ .\ .\ .\ .\ .\ .&\(\varepsilon v\).
\end{tabular}
\end{center}
\par\noindent\hspace*{\parindent}[Setzt man]
\[R+h=r\]
\noindent [so ist:]
\[r^{1-\varepsilon}\cos(1-\varepsilon)v=R^{1-\varepsilon}.\]
\par\noindent\hspace*{\parindent}Also [wird, da \(\delta=(1-\varepsilon)v\) ist:]
\[\left(\tfrac{R+h}{R}\right)^{1-\varepsilon}.\cos\delta=1\]
\noindent [oder]
\[\begin{aligned}
2\sin\tfrac12\delta^2&=1-\left(\tfrac{R}{R+h}\right)^{1-\varepsilon}\\
&=\tfrac{h}{R}(1-\varepsilon)-\tfrac{hh}{RR}\tfrac{(1-\varepsilon)(2-\varepsilon)}{2}\ .\ .\ .\ .
\end{aligned}\]
\par\noindent\hspace*{\parindent}Nimmt man die Bahn des Lichts wie einen Kreisbogen an, so ist
\[\begin{aligned}
\tfrac{R+h}{R}\left[=\tfrac{\sin(v-\delta)}{\sin v-\sin\delta}=\tfrac{\cos\frac12(v-\delta)}{\cos\frac12(v+\delta)}\right]\\
=\frac{\cos\tfrac{\varepsilon}{2-2\varepsilon}.\delta}{\cos\tfrac{2-\varepsilon}{2-2\varepsilon}.\delta}.
\end{aligned}\]
\begin{center}[3.]\end{center}
\par\noindent\hspace*{\parindent}\(A,\ B\)\ .\ .\ .\ .\ Höhen zweier Punkte, wo die gegenseitigen Depressionen
\(a,\ b\) beobachtet werden.
\par\noindent\hspace*{\parindent}\(R\)\ .\ .\ .\ .\ Krümmungshalbmesser [der Erdkugel]; \(r\)\ .\ .\ .\ .\ Entfernung;
\begin{center}\(\mu\)\ .\ .\ .\ .\ Refraktionskoeffizient.\end{center}
\par\noindent\hspace*{\parindent}[Es sei \(v\) der zu der Entfernung \(r\) und \(\gamma\) der zum Lichtstrahl gehörige
Zentriwinkel; dann ist
\endgroup
